\chapter{Section A: Concept Questions}\label{ch:concepts}

Answer briefly. Where a proof or a trace is requested, show it. Model answers are in
Chapter~\ref{ch:solconcepts}; the numbers agree.

\section*{Foundations and array lists}
\begin{questions}
\item\label{q:adt} Define \emph{abstract data type}. State the difference between an ADT and a data
  structure, and name the three ADTs of this coverage that are specialisations or
  relatives of the List.
\item Prove or disprove: (a) $23=O(\log n)$; (b) $5\log n=O(1)$; (c) $5n-7=O(n)$. Give
  constants $c,n_0$ where the claim is true.
\item State the tightest running time, in $\Theta$-notation, of each operation.
  (a)~\cc{push\_back} on an array list that is not full; (b)~insert at the front of an
  array list; (c)~\cc{push\_front} on an SLL; (d)~\cc{pop\_back} on an SLL with head and
  tail pointers; (e)~\cc{pop\_back} on a DLL with head and tail; (f)~worst-case
  \cc{get(i)} on a DLL; (g)~searching an unsorted list; (h)~\cc{push\_front} on a
  circular-array deque that is not full.
\item A function \cc{void f(list l)} receives an array-list struct (with an array inside
  it) by value and sets \cc{l.items[0]=100}. Does the caller see the change? What if
  \cc{l} is a struct with a \cc{head} pointer to nodes, and the function sets
  \cc{l.head->item=100}? What if the parameter is a raw array \cc{int a[]}? Explain each
  answer.
\item After \cc{erase(pos)} on an array list, the old last cell still holds a stale copy of
  the last value. Is this a bug? Which field makes the value harmless, and when would the
  stale value matter?
\item Distinguish \cc{malloc}/\cc{free} from \cc{new}/\cc{delete}. Define a \emph{dangling
  pointer} and a \emph{memory leak} and give a two-line example of each.
\item A function \cc{void append(node *h, int x)} correctly appends to a non-empty list
  but silently fails when the list is empty. Why? Give two different repairs.
\item An array list currently holds $s$ items. How many items are moved by
  \cc{insert(p,x)} (positions $0\le p\le s$) and by \cc{erase(p)}? For which $p$ is each
  worst, and what is that worst case?
\item A dynamic array doubles its capacity when full. Show that $n$ successive
  \cc{push\_back} calls cost $O(n)$ in total. Why does growing by a constant $+10$ instead
  give $\Theta(n^2)$ total?
\item Why is it a mistake to halve the capacity as soon as $n\le\mathit{cap}/2$? Describe
  a bad sequence and a safe shrinking rule.
\item Binary search on a sorted array list needs what precondition and takes what time?
  Explain why a sorted-list insert nevertheless remains $\Theta(n)$ in the worst case
  even though the position is found in $O(\log n)$.
\item In the circular-array queue with fields (\cc{front}, \cc{size}) and capacity
  \cc{CAP}: write the index of the front item, of the back item, and of the next free
  cell. Which values of (\cc{front}, \cc{size}) with $\cc{CAP}=4$ describe a full queue
  whose back item is at index 2?
\end{questions}

\section*{Linked lists}
\begin{questions}[resume]
\item For which operations does adding a tail pointer to a head-only SLL improve the
  running time, and from what to what? Which operation stays $\Theta(n)$ even with a tail
  pointer?
\item Explain precisely why \cc{pop\_back} on an SLL is $\Theta(n)$ even though a tail
  pointer is available.
\item A DLL node \cc{c} must be inserted between neighbouring nodes \cc{a} and \cc{b}
  (with \cc{a->next==b}). List every pointer assignment.
\item List the four cases that occur when removing a node \cc{cur} from a DLL with
  \cc{head}/\cc{tail} pointers, and the assignments for each.
\item What is a sentinel (dummy) node? Name two special cases it removes from list code,
  and one cost.
\item Describe the slow--fast pointer technique and use it to find the middle node of an
  SLL in one pass. What is returned for a list of even length, with the usual loop
  condition?
\item State the loop invariant of the three-pointer in-place reversal of an SLL, and the
  running time and extra space.
\item Define \emph{shallow copy} and \emph{deep copy} of a linked list. If \cc{l = m;} is
  executed for two list structs and \cc{m} is then destroyed node by node, what happens
  to \cc{l}?
\item What is \emph{splicing}? Which pointers change when the whole of list \cc{m} is
  spliced after node \cc{p} of list \cc{l}, and why is the cost $\Theta(1)$ only if some
  pointer is already known?
\item An SLL keeps \cc{head}, \cc{tail} and \cc{mid}. State the invariant for \cc{mid}.
  Which of \cc{push\_front}, \cc{push\_back}, \cc{pop\_front}, \cc{pop\_back},
  \cc{middle} can be $\Theta(1)$ and which not? Which change to the node type would make
  all of them $\Theta(1)$?
\end{questions}

\section*{Stacks}
\begin{questions}[resume]
\item Define the stack ADT and give three applications. Which end of a singly linked list
  should serve as the top, and why?
\item Explain why a queue built from two stacks has amortised $\Theta(1)$ cost per
  operation although a single dequeue may cost $\Theta(n)$.
\item Show with the string \cc{([)]} that counting the opening and closing brackets of
  each kind is not enough to test for balance, and say what data structure is.
\item Convert $(a+b)*c-d/e$ to postfix and to prefix.
\item Evaluate the postfix expression \cc{5 1 2 + 4 * + 3 -}, showing the stack after
  each token.
\item Convert \cc{a-b*c+d/(e-f)} to postfix with the shunting-yard algorithm, showing the
  operator stack after each token. Then explain what changes for \cc{8-3-2} if the
  algorithm pops only operators of \emph{strictly} higher precedence.
\item The numbers $1,2,3,4,5$ are pushed in this order onto a stack, with pops
  interleaved arbitrarily. Which of the following output sequences are possible? For each
  possible one give the push/pop word; for each impossible one name the forbidden pattern.
  (a)~\cc{3 2 5 4 1}; (b)~\cc{4 1 3 2 5}; (c)~\cc{1 2 3 4 5}; (d)~\cc{4 3 5 1 2}. How
  many orders are possible for $1,2,3,4$?
\item Give the next greater element of every item of $[4,5,2,25,7,18]$ (or $-1$), the
  stack after each step of the monotonic-stack algorithm, and an argument that the
  algorithm is $\Theta(n)$ although it contains a loop inside a loop.
\item A min stack stores two stacks \cc{all} and \cc{mins}. Give both after the pushes
  $5,3,7,3,1$, then after two pops, and the value of \cc{min()} at the end. Why is it
  wrong to store only one variable \cc{currentMin}?
\end{questions}

\section*{Queues and deques}
\begin{questions}[resume]
\item A circular queue stores \cc{front} and \cc{rear} (the index \emph{after} the last
  item) instead of \cc{size}, in an array of $\cc{CAP}=6$ cells. (a)~How many items does
  it hold when $\cc{front}=4$, $\cc{rear}=1$? (b)~Why must one cell stay unused, and what
  are the tests for empty and for full?
\item In a linked queue, why is the front the \emph{head} and not the tail? Which field
  must \cc{dequeue} update when it removes the last item, and what happens on the next
  \cc{enqueue} if it does not?
\item Give the values of the C expressions \cc{-7/2}, \cc{-7\%2}, \cc{7\%-2},
  \cc{(0-1)\%5} and \cc{(0-1+5)\%5}. Explain the bug in
  \cc{front = (front - 1) \% CAP;} inside \cc{push\_front}.
\item Starting from an empty deque, perform \cc{push\_front(1)}, \cc{push\_back(2)},
  \cc{push\_front(3)}, \cc{pop\_back()}, \cc{push\_back(4)}, \cc{push\_front(5)},
  \cc{pop\_front()}. Give the value returned by each pop and the final contents from front
  to back.
\item Why can a plain queue not maintain the maximum of a sliding window, while a deque
  can? Give the window maxima of $[9,1,8,2,7,3]$ with $k=2$, and prove that the monotonic
  deque algorithm performs at most $2n$ pushes and pops.
\item An ArrayDeque holds $n=10$ items. How many items are moved by (a)~\cc{add(2,x)},
  (b)~\cc{add(7,x)}, (c)~\cc{remove(5)}, each applied to the original deque? What is the
  worst case, and why is it half that of an array list?
\item A linked deque keeps its sentinel node \emph{inside} the struct
  (\cc{linkedDeque\{ dqnode sentinel; int size; \}}). How many pointer assignments does
  \cc{push\_front} perform on a non-empty deque, and on an empty one? What goes wrong after
  \cc{linkedDeque e = d;}?
\item Seven people numbered $1..7$ stand in a circle and every third person is removed,
  counting from person 1. Give the order of removal and the survivor, and check the
  survivor with the recurrence $J(i)=(J(i-1)+k)\bmod i$.
\item For each task name the most suitable ADT (stack, queue or deque) and say why:
  (a)~undo in an editor; (b)~a printer spooler; (c)~browser history that keeps only the
  last 50 pages but supports \emph{back}; (d)~checking matching HTML tags; (e)~the
  maximum of the last $k$ readings of a sensor; (f)~serving customers in order of arrival.
\end{questions}
